En esta práctica se estudió cómo se dobla un rayo de luz al pasar del aire a una lente de acrílico en forma de media luna. El montaje fue sencillo: una fuente de luz con rendija, una mesa de rayos graduada en grados y la lente colocada en su centro, girando la mesa para cambiar el ángulo con el que llegaba el rayo. Como la sesión de laboratorio no pudo realizarse por falta de equipo y de profesor, los ángulos de refracción que aparecen en este reporte son datos simulados, calculados con la ley de Snell para el acrílico y modificados con una desviación de entre 5 y 10 por ciento para imitar una medición real. Con ellos se graficó el seno del ángulo de incidencia contra el seno del ángulo de refracción. La recta obtenida tiene una pendiente de 1.396, que es el índice de refracción del acrílico, frente al valor de referencia de 1.49. El error es de 6.3 por ciento.
